% Generated draft: requires mathematical and source review.
\documentclass[11pt]{article}
\usepackage[margin=25mm]{geometry}
\usepackage{fontspec}
\setmainfont{Latin Modern Roman}
\usepackage{amsmath,amssymb,mathtools,unicode-math}
\setmathfont{Latin Modern Math}
\usepackage{enumitem,needspace,xurl,hyperref}
\setlength{\parindent}{0pt}
\setlength{\parskip}{5pt}
\setlength{\emergencystretch}{3em}
\newcommand{\N}{\mathbb{N}}
\newcommand{\Z}{\mathbb{Z}}
\newcommand{\Q}{\mathbb{Q}}
\newcommand{\R}{\mathbb{R}}
\newcommand{\C}{\mathbb{C}}
\newcommand{\F}{\mathbb{F}}
\newcommand{\A}{\mathbb{A}}
\newcommand{\PP}{\mathbb{P}}
\newcommand{\E}{\mathbb{E}}
\newcommand{\eps}{\varepsilon}
\newcommand{\dd}{\,\mathrm{d}}
\begin{document}
\section*{UnsolvedMath ID: 30001469}
\textbf{Source problem code:} catalogue entry
\subsection{Definitions and mathematical statement}
Let \(G\) be a group and let \(kG\) denote its group ring over the coefficient ring \(k\). Regard \(k\) as the trivial \(kG\)-module via the augmentation. Write
\[
\operatorname{cd}_k(G)=\operatorname{pd}_{kG}(k),\qquad
\operatorname{hd}_k(G)=\operatorname{fd}_{kG}(k),
\]
where projective and flat dimensions may be infinite. For a finite subgroup \(F\leq G\), invertibility of its order means \(|F|1_k\in k^\times\).

An elementary amenable group belongs to the smallest class containing finite groups and abelian groups and closed under extensions and directed unions. Write \(h(G)\) for its extended Hirsch length: this agrees with \(\dim_{\Q}(\Q\otimes_{\Z}A)\) for an abelian group \(A\), is additive under extensions, and is the supremum of the Hirsch lengths in a directed union. A group is locally finite if each of its finitely generated subgroups is finite. Write \(\aleph_\omega=\sup_{n<\omega}\aleph_n\). A filtered colimit below is a colimit in the category of groups of a diagram indexed by a small filtered category; the transition homomorphisms are not assumed injective.

\textbf{Coefficient qualification.} The supplied original statement specifies a field \(k\), whereas the supplied statement and clean statement specify a commutative ring \(k\). The following records the proposed formulas with this discrepancy unresolved; the commutative-ring version must not be identified silently with the field version. The ring conventions concerning an identity and exclusion of the zero ring are not supplied.

\textbf{First proposed equivalence.} For every elementary amenable group \(G\), is
\[
\operatorname{cd}_k(G)<\infty
\quad\Longleftrightarrow\quad
\left[
 h(G)<\infty,\quad |G|<\aleph_\omega,\quad
 |F|1_k\in k^\times\text{ for every finite }F\leq G
\right]?
\]

\textbf{Second proposed equivalence.} For every group \(G\), without the elementary amenability assumption, is
\[
\operatorname{hd}_k(G)\leq 1
\quad\Longleftrightarrow\quad
\left[
 G\cong\varinjlim_{i\in I}G_i
 \text{ for a small filtered diagram with }
 \operatorname{cd}_k(G_i)\leq 1\text{ for every }i
\right]?
\]

\textbf{Third proposed equivalence.} For an elementary amenable group \(G\), let \(\mathcal A_k(G)\) denote the supplied action condition: \(G\) admits an action on a contractible complex of dimension \(2\), with finite stabilizers whose orders are invertible in \(k\). The category of complexes, regularity of the action, and meaning of stabilizer are not specified in the record.

Define the three listed structural conditions as follows:
\begin{enumerate}
\item \(C_1(G)\): there exist a subgroup \(H\leq G\) with \([G:H]<\infty\) and an integer \(n\in\Z\setminus\{0\}\) such that
\[
H\cong\langle x,y\mid y^{-1}xy=x^n\rangle.
\]
\item \(C_2(G)\): there exists a countably infinite locally finite normal subgroup \(T\trianglelefteq G\) such that \(G/T\cong\Z\) or \(G/T\cong D_\infty\), where \(D_\infty=\Z\rtimes(\Z/2\Z)\) and the nonidentity element of \(\Z/2\Z\) acts on \(\Z\) by inversion.
\item \(C_3(G)\): \(G\) is locally finite and \(|G|=\aleph_1\).
\end{enumerate}
The clean statement asks whether
\[
\operatorname{cd}_k(G)=2
\quad\Longleftrightarrow\quad\mathcal A_k(G)
\quad\Longleftrightarrow\quad
C_1(G)\lor C_2(G)\lor C_3(G).
\]
The original statement instead says that \emph{exactly one} of \(C_1(G),C_2(G),C_3(G)\) holds. This exclusivity qualification is retained as an unresolved difference between the supplied formulations.

\textbf{Further unresolved qualifications.} No minimality condition on the dimension of the contractible complex is supplied, although the asserted cohomological dimension is exactly \(2\). Also, the structural alternatives do not explicitly impose invertibility in \(k\) of every finite-subgroup order, while that restriction occurs elsewhere in the record. It is unresolved whether a global restriction of this kind, or another omitted hypothesis, was intended for the third equivalence. These are proposed source formulations requiring review, not asserted theorems.
\subsection{Short English statement}
The record asks for three characterizations: finite cohomological dimension for elementary amenable groups, homological dimension at most one for arbitrary groups, and cohomological dimension exactly two for elementary amenable groups. The last characterization uses both a contractible complex action and three structural alternatives. The coefficient domain, exclusivity of those alternatives, and qualifications on the action and torsion remain unresolved. The supplied status is \emph{partially solved}; this draft does not independently verify that assessment.
\subsection{Sources}
\begin{itemize}
\item Supplied local UnsolvedMath registry record: permanent ID \(30001469\), title \emph{Cohomological Dimensions of Elementary Amenable Groups}, exact source problem code \texttt{catalogue entry}, in the \emph{Oberwolfach Reports Open Problems} collection. The record gives proposed year \(2010\) and names no individual proposer.
\item Supplied report citation: \emph{Non-positive Curvature and Geometric Structures in Group Theory} (2010), Oberwolfach Reports, pp.\ 1165--1224, DOI \texttt{10.4171/owr/2010/20}. Supplied source URL and literature-source link labelled ``original OWR report'': \url{https://doi.org/10.4171/owr/2010/20}. The background describes this report as documenting workshop discussions and open problems.
\item Supplied literature-source link labelled ``Bridson--Kropholler, dimension'': \url{https://arxiv.org/abs/1208.1084}.
\item Literature checked at: \texttt{2026-08-21}. This is supplied dated curation, not an independently verified literature assessment. The supplied assessment labels the bundled problem \emph{partially solved}, attributes substantial special cases to Bridson--Kropholler, including elementary amenable groups of homological dimension one, and describes the general finite cohomological-dimension classification as remaining conjectural. The background explicitly calls this dated literature triage and requests independent verification of time-sensitive status and cited claims.
\item The supplied statement-status metadata is \texttt{corrected\_verified}; its verification note reports checking against the official Oberwolfach report and preserving the hypotheses, notation, and open target. This is a supplied verification claim, not verification performed for this draft. The record was updated at \texttt{2026-09-16T12:06:53Z}.
\end{itemize}
\end{document}
